\directlua{dofile('processing.lua')}

\section{Прямые однократные измерения}

Абсолютная погрешность измерения напряжения $\Delta U$ определяется приборной погрешностью вольтметра, заданной в исходных данных:
\begin{equation}
    \Delta U = \Delta_V = \directlua{tex.print(Delta_instr)} \text{ В}
\end{equation}

Относительная погрешность измерения напряжения $\delta_U$ рассчитывается по формуле:
\begin{equation}
    \delta_U = \frac{\Delta U}{U} \cdot 100\% = \frac{\directlua{tex.print(Delta_instr)}}{\directlua{tex.print(U_s)}} \cdot 100\% = \directlua{tex.print(round(rel_delta_U, 4))} \%
\end{equation}

Результат однократного измерения напряжения:
\begin{equation*}
    { U_x = (\directlua{tex.print(U_s)} \pm \directlua{tex.print(Delta_instr)}) \text{ В}, \quad \delta_U = \directlua{tex.print(round(rel_delta_U, 3))} \% }
\end{equation*}

\section{Косвенные однократные измерения}

\subsection{Измерение тока}
Ток $I$, протекающий через делитель, рассчитывается по падению напряжения $U_0$ на образцовом резисторе $R_0 = \directlua{tex.print(R0)}$ Ом.

\begin{equation}
    I = \frac{U_0}{R_0} = \frac{\directlua{tex.print(U0_s)}}{\directlua{tex.print(R0)}} = \directlua{tex.print(round(I_s, 5))} \text{ А}
\end{equation}

Абсолютная погрешность измерения напряжения на образцовом резисторе равна приборной погрешности: $\Delta U_0 = \directlua{tex.print(Delta_instr)}$ В.
Относительная погрешность измерения напряжения $\delta_{U_0}$ согласно документации вольтметра $\delta_{U_o}=c + d\cdot\left[ \frac{X_k}{X} - 1 \right]$, где $X_k$ --- пределы измерения; $X$ --- измеренная величина; $c$ и $d$ -- класс точности: 
\begin{equation}
    \delta_{U_0} = \directlua{tex.print(ktc)} + \directlua{tex.print(ktd)}\cdot\left[ \frac{\directlua{tex.print(X_k)}}{\directlua{tex.print(U0_s)}} - 1 \right] = \directlua{tex.print(round(rel_delta_U0, 2))} \%
\end{equation}

Относительная погрешность косвенного измерения тока $\delta_I$ (сумма относительных погрешностей аргументов):
\begin{equation}
    \delta_I = \delta_{U_0} + \delta_{R_0} = \directlua{tex.print(round(rel_delta_U0, 2))} + \directlua{tex.print(delta_R0)} = \directlua{tex.print(round(delta_I_rel, 2))} \%
\end{equation}

Абсолютная погрешность тока $\Delta I$:
\begin{equation}
    \Delta I = I \cdot \frac{\delta_I}{100} = \directlua{tex.print(round(I_s, 5))} \cdot \frac{\directlua{tex.print(round(delta_I_rel, 3))}}{100} = 
    \luaexec{tex.print(string.format('\%.8f', abs_delta_I))}
    \text{ А}
\end{equation}

Результат измерения тока:
\begin{equation*}
    { I_x = (\directlua{tex.print(round(I_s, 5)*1000)} \pm \luaexec{tex.print(string.format('\%.2f', round(abs_delta_I, 5)*1000))}) \cdot 10^3\text{ А} }
\end{equation*}

\subsection{Измерение мощности}
Мощность $P_M$, выделяемая на участке делителя:
\begin{equation}
    P_M = U \cdot I = \directlua{tex.print(U_s)} \cdot \directlua{tex.print(round(I_s, 5))} = \directlua{tex.print(round(P_s, 4))} \text{ Вт}
\end{equation}

Относительная погрешность мощности $\delta_P$:
\begin{equation}
    \delta_P = \delta_U + \delta_I = \directlua{tex.print(round(rel_delta_U, 3))} + \directlua{tex.print(round(delta_I_rel, 3))} = \directlua{tex.print(round(delta_P_rel, 3))} \%
\end{equation}

Абсолютная погрешность $\Delta P_M$:
\begin{equation}
    \Delta P_M = P_M \cdot \frac{\delta_P}{100} = \directlua{tex.print(round(math.abs(P_s), 4))} \cdot \frac{\directlua{tex.print(round(delta_P_rel, 3))}}{100} = \directlua{tex.print(round(abs_delta_P, 4))} \text{ Вт}
\end{equation}

Результат измерения мощности:
\begin{equation*}
    { P_{Mx} = (\directlua{tex.print(round(P_s, 3))} \pm \directlua{tex.print(round(abs_delta_P, 3))}) \text{ Вт} }
\end{equation*}

\section{Обработка многократных измерений (Напряжение)}

Число наблюдений $n = \directlua{tex.print(n)}$. Доверительная вероятность $P = \directlua{tex.print(prob_P)}$. Коэффициент Стьюдента $t_p(n-1) = \frac{t_{95\%}(8) + t_{95\%}(10)}{2} = \frac{\directlua{tex.print(t_8)} + \directlua{tex.print(t_10)}}{2} = \directlua{tex.print(t_val)}$.

Оценка математического ожидания (среднее арифметическое):
\begin{equation}
    \bar{U} = \frac{1}{n}\sum_{i=1}^{n} U_{1i} = \directlua{tex.print(round(U_mean, 4))} \text{ В}
\end{equation}

Оценка среднеквадратического отклонения (СКО) результата наблюдений $S[U]$ и среднего арифметического $S[\bar{U}]$:
\begin{equation}
    S[U] = \sqrt{\frac{\sum (U_{1i} - \bar{U})^2}{n-1}} = \directlua{tex.print(round(S_U, 4))} \text{ В}
\end{equation}
\begin{equation}
    S[\bar{U}] = \frac{S[U]}{\sqrt{n}} = \frac{\directlua{tex.print(round(S_U, 4))}}{\sqrt{\directlua{tex.print(n)}}} = \directlua{tex.print(round(S_U_mean, 4))} \text{ В}
\end{equation}

\begin{luacode*}
    function t89()
        for i = 1, #U1_arr do
            local ui = U1_arr[i]
            local diff = ui - U_mean
            local diff_sq = diff * diff
            tex.print(string.format("$%d$ & $% .6f$ & $%.6f$ & $%.6f$ \\\\", i, ui, diff, diff_sq))
        end
    end
\end{luacode*}

\begin{table}[H]
    \centering
    \caption{Результаты измерений и вычислений для напряжения}
    \label{tab:5_1}
    \begin{tblr}[expand=\directlua]{
        colspec = { Q[c,m, 3em] X[c,m] X[c,m] X[c,m] },
        hlines, vlines,
    }
        Номер измерения & $U_{1i}, \text{В}$ & $(U_{1i} - \bar{U}), \text{В}$ & $(U_{1i} - \bar{U})^2, \text{В}^2$ \\
        \directlua{t89()}
    \end{tblr}
\end{table}

Доверительный интервал $\Delta U$ (случайная составляющая):
\begin{equation}
    \Delta U = t_p(f) \cdot S[\bar{U}] = \directlua{tex.print(t_val)} \cdot \directlua{tex.print(round(S_U_mean, 4))} = \directlua{tex.print(round(delta_U_multi, 4))} \text{ В}
\end{equation}

Результат многократных измерений напряжения:
\begin{equation*}
    { U_x = (\directlua{tex.print(round(U_mean, 2))} \pm \directlua{tex.print(round(delta_U_multi, 2))}) \text{ В}, \quad P = \directlua{tex.print(prob_P)} }
\end{equation*}

\section{Многократные косвенные измерения}

\subsection{Ток}
Массив значений тока получен как $I_i = U_{2i} / R_0$ (\cref{tab:current_massive}).

\begin{luacode*}
    function t104()
        local n = #U2_arr
        local half = math.ceil(n / 2)
        for i = 1, half do
            local i2 = i + half
            local val1 = string.format("%.5f", U2_arr[i] / R0)
            
            -- Проверка существования элемента для второй колонки
            local idx2 = ""
            local val2 = ""
            if i2 <= n then
                idx2 = tostring(i2)
                val2 = string.format("%.5f", U2_arr[i2] / R0)
            end
            
            tex.print(string.format("$%d$ & $%s$ & $%s$ & $%s$ \\\\", i, val1, idx2, val2))
            
        end
    end
\end{luacode*}
 
\begin{table}[H]
    \centering
    \caption{Массив расчетных значений тока $I_i$, А}
    \label{tab:current_massive}
    \begin{tblr}[expand=\directlua]{
        colspec = { Q[c,m, 2em] X[c,m] Q[c,m, 2em] X[c,m] },
        hlines, vlines,
    }
        № & $I_i$ & № & $I_i$ \\
        \directlua{t104()}
    \end{tblr}
\end{table}

Среднее значение тока:
\begin{equation}
    \bar{I} = \frac{1}{n}\sum_{i=1}^{n} I_i = \directlua{tex.print(round(I_mean, 5))} \text{ А}
\end{equation}

\begin{luacode*}
    function t168()
        for i = 1, #U1_arr do
            local u1 = U1_arr[i]
            local u2 = U2_arr[i]
            local diff_u = u1 - U_mean
            local cur_i = u2 / R0
            local diff_i = cur_i - I_mean
            tex.print(string.format("$%d$ & $%.3f$ & $%.3f$ & $%.3f$ & $%.5f$ & $%.5f$
             \\\\", 
                i, u1, diff_u, u2, cur_i, diff_i))
        end
    end
\end{luacode*}

\begin{table}[H]
    \centering
    \caption{Результаты измерений и вычислений для тока}
    \label{tab:5_3}
    \begin{tblr}[expand=\directlua]{
        colspec = { Q[c,m, 2em] *{5}{X[c,m]} },
        hlines, vlines,
        rows = {font=\small},
    }
        Номер & $U_{1i}, \text{В}$ & $(U_{1i} - \bar{U}), \text{В}$ & $U_{2i}, \text{В}$ & $I_i, \text{А}$ & $(I_i - \bar{I}), \text{А}$ \\
        \directlua{t168()}
    \end{tblr}
\end{table}

СКО среднего арифметического тока:
\begin{equation}
    S[\bar{I}] = \frac{S[I]}{\sqrt{n}} = \frac{\directlua{tex.print(round(S_I, 6))}}{\directlua{tex.print(round(math.sqrt(n), 6))}} = \directlua{tex.print(round(S_I_mean, 6))} \text{ А}
\end{equation}

Доверительный интервал для тока ($t_p = \directlua{tex.print(t_val)}$):
\begin{equation}
    \Delta I = t_p \cdot S[\bar{I}] = \directlua{tex.print(t_val)} \cdot \directlua{tex.print(round(S_I_mean, 6))} = \directlua{tex.print(round(delta_I_multi, 5))} \text{ А}
\end{equation}

Результат:
\begin{equation*}
    { I_x = (\directlua{tex.print(round(I_mean, 3))} \pm \directlua{tex.print(round(delta_I_multi, 3))}) \text{ А} }
\end{equation*}

\subsection{Мощность}
Средняя мощность:
\begin{equation}
    P_M = \bar{U} \cdot \bar{I} = \directlua{tex.print(round(U_mean, 4))} \cdot \directlua{tex.print(round(I_mean, 5))} = \directlua{tex.print(round(P_multi_mean, 4))} \text{ Вт}
\end{equation}

Оценка СКО косвенного измерения мощности:
\begin{equation}
    S[P_M] = \sqrt{\bar{U}^2 S^2[\bar{I}] + \bar{I}^2 S^2[\bar{U}]} = \directlua{tex.print(round(S_Pm, 5))} \text{ Вт}
\end{equation}

Эффективное число степеней свободы $f_{\text{эф}}$:
\begin{equation}
    f_{\text{эф}} = (n+1)\frac{(\bar{U}^2 S^2[\bar{I}] + \bar{I}^2 S^2[\bar{U}])^2}{\bar{U}^4 S^4[\bar{I}] + \bar{I}^4 S^4[\bar{U}]} - 2 = \directlua{tex.print(math.floor(f_eff))}
\end{equation}

Коэффициент Стьюдента для $f_{\text{эф}} = \directlua{tex.print(math.floor(f_eff))}$ равен $k_p = \directlua{tex.print(t_val_eff)}$.
Доверительный интервал:
\begin{equation}
    \Delta P_M = k_p \cdot S[P_M] = \directlua{tex.print(t_val_eff)} \cdot \directlua{tex.print(round(S_Pm, 5))} = \directlua{tex.print(round(delta_P_multi, 4))} \text{ Вт}
\end{equation}

Результат:
\begin{equation*}
    { P_{Mx} = (\directlua{tex.print(round(P_multi_mean, 2))} \pm \directlua{tex.print(round(delta_P_multi, 2))}) \text{ Вт}, \quad P = \directlua{tex.print(prob_P)} }
\end{equation*}

\section*{Выводы}

\begin{enumerate}
    \item В ходе прямых и косвенных однократных измерений установлено влияние отношения измеряемой величины к точноти измерений прибора. Относительная погрешность прямого измерения напряжения составила $\delta_U = \directlua{tex.print(round(rel_delta_U, 3))}\%$, при этом погрешность косвенного измерения тока достигла $\delta_I = \directlua{tex.print(round(delta_I_rel, 3))}\%$. Данное расхождение вызвано тем, что измеряемое падение напряжения на образцовом резисторе ($U_0 = \directlua{tex.print(U0_s)} \text{ В}$) соизмеримо с абсолютной инструментальной погрешностью вольтметра ($\Delta_V = \directlua{tex.print(abs_delta_U0)} \text{ В}$).
    
    \item При обработке многократных измерений установлено, что наличие случайных погрешностей существенно увеличивает разброс результатов. При сильном зашумлении сигнала (где относительная погрешность измерения тока достигает $\delta_I = \directlua{tex.print(round(delta_I_rel, 2))}\%$) использование статистических методов обработки позволяет корректно определить доверительный интервал для истинного значения измеряемой величины с заданной вероятностью $P = 0.95 \%$.
    
    \item При косвенных многократных измерениях мощности расчет эффективного числа степеней свободы ($f_{\text{эф}} = \directlua{tex.print(round(f_eff, 0))}$) подтвердил, что итоговая дисперсия результата формируется аргументами неравномерно.
\end{enumerate}